% Standalone proof. Shared notation and prose conventions: analyses/README.md.
\documentclass[11pt,a4paper]{article}
\usepackage[T1]{fontenc}
\usepackage[utf8]{inputenc}
\usepackage{amsmath,amsthm}
\usepackage{newtxtext,newtxmath}
\usepackage[a4paper,margin=21mm,headheight=15pt]{geometry}
\usepackage{microtype,booktabs,array,tabularx,longtable}
\usepackage[table]{xcolor}
\usepackage{tikz,fancyhdr,enumitem,seqsplit}
\definecolor{Teal}{HTML}{176B70}
\definecolor{Ink}{HTML}{203443}
\definecolor{Pale}{HTML}{EDF5F5}
\definecolor{Line}{HTML}{C7D8DB}
\definecolor{Muted}{HTML}{536974}
\definecolor{Warm}{HTML}{F8F4EC}
\usepackage[colorlinks=true,linkcolor=Teal,urlcolor=Teal]{hyperref}
\color{Ink}
\newcommand{\blank}{\lambda}
\newcommand{\digits}{\{0,1\}}
\newcommand{\ph}[1]{\mathsf{#1}}
\newcommand{\EXPERT}{\mathcal E}
\newcommand{\CTRL}{\mathcal C}
\newcommand{\PTO}{\operatorname{PTO}}
\newcommand{\triple}[3]{(#1,#2,\mathsf{#3})}
\newcommand{\code}[1]{\texttt{#1}}
\newcommand{\hash}[1]{{\scriptsize\ttfamily\seqsplit{#1}}}
\newtheorem{theorem}{Theorem}
\newtheorem{lemma}{Lemma}
\newtheorem{corollary}{Corollary}
\setlength{\parindent}{0pt}
\setlength{\parskip}{6pt}
\setlist[itemize]{nosep,leftmargin=1.2em}
\pagestyle{fancy}
\fancyhf{}
\fancyhead[L]{\small\sffamily\color{Muted}Binary Addition}
\fancyhead[R]{\small\sffamily\color{Muted}Metric Extraction}
\fancyfoot[L]{\scriptsize\sffamily\color{Muted}Exact PTO Equivalence}
\fancyfoot[R]{\small\thepage}
\renewcommand{\headrulewidth}{.35pt}
\renewcommand{\footrulewidth}{0pt}
\newcommand{\sectiontitle}[2]{\refstepcounter{section}\vspace*{2pt}{\LARGE\bfseries\color{Teal}\thesection\quad #2}\par\vspace{7pt}}
\date{}
\pdfinfoomitdate=1
\pdftrailerid{}
\begin{document}
{\Huge\bfseries Exact PTO Equivalence\\[3pt]Binary Addition}\par
\vspace{10pt}
{\Large Metric-Extracted Controller}\par
{\small\color{Muted}\path{analyses/metric_extraction/DFST/add2/warmup_50000_steps_250000/}}\par
\vspace{7pt}

\vspace{4pt}
\sectiontitle{1 \enspace STATEMENT AND SCOPE}{Statement and Scope}\label{sec:common1}
\textbf{Controllers.}
The exported table \path{minimized_transition_table.pdf} defines controller
$\CTRL$, with initial state $0$ and halt state $21$.
The function \code{add\_grid(2)} defines expert $\EXPERT$, with initial
state $0$ and halt state $10$. The adjacent expert export matches
this function entry for entry. All tables retain the exported state numbers.

\textbf{Input Domain.}
Let $D=\digits$. Row $0$ initially contains the contiguous word $a+b$,
starting at column $0$, with $a\in D^m$, $b\in D^n$, and $m,n\ge1$.
All remaining cells contain blanks. The head starts at $(0,0)$.
Write $A=\operatorname{val}_{2}(a)$ and $B=\operatorname{val}_{2}(b)$.
The domain permits arbitrary operand lengths, unequal lengths, and leading zeros.

\textbf{Execution Convention.}
The grid uses $\mathbb Z^2$, with coordinates increasing rightward and
downward. The symbol $\blank$ denotes a blank. A transition
$(q',w,M)$ writes $w$, performs motion $M$, and enters state $q'$.
Motions $L,R,D,U$ move one cell. Motion $S$ writes and stops without moving.

\textbf{Policy-Trajectory Observations.}
For input $a+b$, define the policy-trajectory observation (PTO) by
\[
\PTO(a+b)=\bigl((r_t,w_t,M_t)\bigr)_{t=0}^{T-1},
\]
where $r_t$, $w_t$, and $M_t$ denote the read symbol, written symbol,
and motion at controller step $t$. A finite execution contributes the final
$S$ action exactly once, without halt-loop padding.

\begin{theorem}[Exact PTO Equivalence]\label{thm:main}
For every input in the stated domain,
$\PTO_{\CTRL}(a+b)=\PTO_{\EXPERT}(a+b)$.
Both executions terminate on the same step. The complete grids and head
positions agree at every step, including the final configuration.
\end{theorem}

The theorem concerns autonomous executions of the exported discrete controller
on the stated input domain.
\clearpage
\sectiontitle{2 \enspace EXPERT BEHAVIOR}{Expert Behavior}\label{sec:common2}
\subsection{Accumulator and Marker Invariant}\label{sec:part1}
Write $B=\sum_{j=0}^{n-1}b_j2^j$, with $b_0$ the rightmost digit.
The expert scans to the separator at column $m$, steps left and down,
writes a working $+$ below column $m-1$, and enters state $8$ above the marker.

\begin{lemma}[Round Invariant]\label{lem:invariant}
After processing $k$ right-operand digits, $0\le k\le n$, the expert enters
state $8$ at $(p_k,0)$, where $p_k=m-1-k$. At this point:
\begin{itemize}
\item Row $1$ contains exactly one nonblank symbol: a $+$ at column $p_k$.
\item Row $0$ contains a contiguous accumulator on $[\ell_k,m-1]$, with value
$C_k=A+\sum_{j<k}b_j2^j$ and left edge $\ell_k\le p_k+1$.
Column $m-1$ carries the units digit.
\item The separator remains at column $m$. The unprocessed prefix of $b$
occupies columns $m+1,\ldots,m+n-k$. All remaining cells contain blanks.
\end{itemize}
\end{lemma}
\begin{proof}
Initialization gives $k=0$, $\ell_0=0$, $C_0=A$, and $p_0=m-1$.
State $8$ preserves the digit at $p_k$, or writes $0$ on a blank.
A blank at $p_k$ forces $\ell_k=p_k+1$, so zero extension preserves contiguity.

For $k<n$, the right scan reaches the blank after the unprocessed prefix,
steps left to digit $b_k$ at $r_k=m+n-k$, erases $b_k$, and descends.
The expert transports $b_k$ left along row $1$ to the marker. The distance satisfies
\[
r_k-p_k=n+1\ge2. \tag{1}
\]
Consequently, the first two transport steps read blanks. At the marker, the
expert ascends and adds $b_k$ to the accumulator digit $e$ at column $p_k$.
If $e+b_k<2$, the expert writes $e+b_k$ and descends. Otherwise, the expert
writes $e+b_k-2$, moves left, and propagates one carry. Each digit $1$
changes to $0$; the first smaller digit increases by one, or a blank changes
to $1$. The update adds exactly $b_k2^k$ and preserves contiguity.

The expert returns along row $1$, erases the marker, steps left, writes the
next marker at $p_{k+1}=p_k-1$, and ascends into state $8$.
The earlier zero extension gives $\ell_{k+1}\le p_k=p_{k+1}+1$.
The accumulator now has value $C_{k+1}$, establishing the next round.
\end{proof}

\begin{lemma}[Termination and Final Configuration]\label{lem:termination}
The expert terminates with only the padded base-$2$ representation of $A+B$
on row $0$, ending at column $m-1$. The final head stops at $(m-n-1,0)$.
Writing $\ell_2(s)$ for the binary digit count, with $\ell_2(0)=1$,
the output width equals $L=\max\{m,n+1,\ell_2(A+B)\}$.
\end{lemma}
\begin{proof}
Each round consumes one of the $n$ right-operand digits. Every scan crosses
a finite interval, and every carry reaches a digit below $1$ or a blank.
After round $n$, state $8$ ensures a digit at $p_n$. The next right scan
returns to the separator. State $3$ erases the separator and descends.
State $9$ walks left, erases the final marker, ascends, preserves the
accumulator digit, and stops. The invariant gives $C_n=A+B$, and cleanup
removes all work symbols.
The original $m$ columns persist, and marker extension reaches $p_n$,
ensuring at least $n+1$ digits. Every additional carry extension contributes
a digit required by the current numerical value. Partial sums never decrease,
so the final width equals $L$.
\end{proof}

\clearpage
\sectiontitle{3 \enspace OBSERVATION PROTOCOL}{Observation Protocol}\label{sec:common3}
\label{sec:protocol}
The following phases refine the expert states by geometric context.
The protocol covers every valid execution and may also admit additional
observation sequences. Write $d\in\digits$. An edge $x:j$ reads $x$ and
advances to phase $j$; $d:j$ permits either digit. Expert state $g_i$
supplies the write and motion from Section~\ref{sec:tables}.
A dash denotes a terminal phase with no outgoing transition.

\begin{center}
\small\renewcommand{\arraystretch}{1.21}
\begin{tabular}{r l c l}
\toprule
\rowcolor{Pale}$i$ & Phase meaning & $g_i$ & Allowed reading $:$ next phase\\
\midrule
0 & First input digit & 0 & $d:1$ \\
1 & Seek the separator & 0 & $d:1\quad +:2$ \\
2 & Initial descent & 1 & $d:3$ \\
3 & Place the marker & 1 & $\blank:4$ \\
4 & Ensure an accumulator digit & 8 & $d:6\quad \blank:5$ \\
5 & Digit after zero extension & 2 & $d:6$ \\
6 & Scan toward separator & 2 & $d:6\quad +:7$ \\
7 & First cell after separator & 2 & $d:8\quad \blank:22$ \\
8 & Scan remaining right operand & 2 & $d:8\quad \blank:9$ \\
9 & Read rightmost factor digit & 3 & $0:10\quad 1:14$ \\
10 & Transport zero: first blank & 4 & $\blank:11$ \\
11 & Transport zero: second blank & 4 & $\blank:12$ \\
12 & Transport zero: walk left & 4 & $\blank:12\quad +:13$ \\
13 & Add zero above marker & 4 & $d:21$ \\
14 & Transport one: first blank & 5 & $\blank:15$ \\
15 & Transport one: second blank & 5 & $\blank:16$ \\
16 & Transport one: walk left & 5 & $\blank:16\quad +:17$ \\
17 & Add one above marker & 5 & $0:21\quad 1:18$ \\
18 & Propagate a carry & 6 & $1:18\quad 0:19\quad \blank:19$ \\
19 & First return cell after carry & 7 & $\blank:20$ \\
20 & Return to the marker & 7 & $\blank:20\quad +:3$ \\
21 & At marker, without carry & 7 & $+:3$ \\
22 & Erase the input separator & 3 & $+:23$ \\
23 & First cleanup cell & 9 & $\blank:24$ \\
24 & Find and erase final marker & 9 & $\blank:24\quad +:25$ \\
25 & Preserve result and stop & 9 & $d:26$ \\
26 & Halted & 10 & $\text{--}$ \\
\bottomrule
\end{tabular}
\end{center}

\begin{lemma}[Protocol Coverage]\label{lem:protocol}
Every execution of $\EXPERT$ from a valid input follows this protocol, starting
at phase $0$ and ending at phase $26$.
\end{lemma}
\begin{proof}
The nonempty first operand establishes phases $0$--$3$.
Lemma~\ref{lem:invariant} places a digit immediately after a zero extension
in phase $4$, establishing phase $5$. The right scan meets the separator before any blank.
To the separator's right, a first blank means the next left step reads the separator
(phase $22$); after any digit, that step reads a digit (phase $9$).
Equation~(1) guarantees two initial blank readings while transporting either
bit. The cell above the marker contains a digit. A carry ends strictly left
of the marker, so the first return cell contains a blank. A noncarry addition returns directly to
the marker. Erasing and shifting the marker reaches a blank. Finally,
cleanup starts below the separator on a blank, reaches the marker and then
a digit. These facts justify every branch of the table. Induction over the
steps and rounds proves coverage; Lemma~\ref{lem:termination} supplies the final phase.
\end{proof}

\clearpage
\sectiontitle{4 \enspace SIMULATION CERTIFICATE}{Simulation Certificate}\label{sec:common4}
\subsection{Phase Relation and Local Checks}\label{sec:part3}
For phase $i$, write $Q_i$ for the controller-state set in the following table.
The relation pairs an expert configuration in phase $i$ with a controller
configuration having the \emph{same grid and head} and a state in $Q_i$.

\begin{center}
\small\renewcommand{\arraystretch}{1.18}
\begin{tabular}{r l c l}
\toprule
\rowcolor{Pale}Phase $i$ & Meaning & Expert state $g_i$ & Controller states $Q_i$\\
\midrule
0 & First input digit & 0 & $\{0\}$ \\
1 & Seek the separator & 0 & $\{1\}$ \\
2 & Initial descent & 1 & $\{2\}$ \\
3 & Place the marker & 1 & $\{3\}$ \\
4 & Ensure an accumulator digit & 8 & $\{4\}$ \\
5 & Digit after zero extension & 2 & $\{15\}$ \\
6 & Scan toward separator & 2 & $\{5\}$ \\
7 & First cell after separator & 2 & $\{6\}$ \\
8 & Scan remaining right operand & 2 & $\{6\}$ \\
9 & Read rightmost factor digit & 3 & $\{7\}$ \\
10 & Transport zero: first blank & 4 & $\{12\}$ \\
11 & Transport zero: second blank & 4 & $\{13\}$ \\
12 & Transport zero: walk left & 4 & $\{13\}$ \\
13 & Add zero above marker & 4 & $\{14\}$ \\
14 & Transport one: first blank & 5 & $\{8\}$ \\
15 & Transport one: second blank & 5 & $\{9\}$ \\
16 & Transport one: walk left & 5 & $\{9\}$ \\
17 & Add one above marker & 5 & $\{10\}$ \\
18 & Propagate a carry & 6 & $\{18\}$ \\
19 & First return cell after carry & 7 & $\{19,20\}$ \\
20 & Return to the marker & 7 & $\{20\}$ \\
21 & At marker, without carry & 7 & $\{11\}$ \\
22 & Erase the input separator & 3 & $\{7\}$ \\
23 & First cleanup cell & 9 & $\{16\}$ \\
24 & Find and erase final marker & 9 & $\{16\}$ \\
25 & Preserve result and stop & 9 & $\{17\}$ \\
26 & Halted & 10 & $\{21\}$ \\
\bottomrule
\end{tabular}
\end{center}

\textbf{Local Transition Identity.} For every listed $q\in Q_i$ and every allowed
edge $i\xrightarrow{x}j$ from the protocol, the actual tables satisfy
\begin{equation}
\delta_\EXPERT(g_i,x)=(g_j,w,M),\qquad
\delta_\CTRL(q,x)=(q',w,M),\qquad q'\in Q_j. \tag{2}
\end{equation}
Every required controller entry exists. A transition reaches the halt
state exactly when $j=26$, with motion $S$.
The certificate contains \textbf{28 phase--state pairs} (including the terminal pair)
and \textbf{50 local checks}. The checks cover 43 distinct nonterminal
controller entries. Section~\ref{sec:tables} provides the tables for direct
inspection; the accompanying script checks every instance of~(2).

\newpage
\subsection{Transition Tables}\label{sec:tables}
Each entry lists $(\text{next state},\text{written symbol},\text{motion})$.
The symbol $\blank$ means blank; a dash means an undefined transition.
The tables reproduce the adjacent exports with the original state numbers.

\textbf{Expert: \code{add\_grid(2)}.}\quad Initial state $0$; halt $10$.
\begin{center}
\fontsize{9.4}{11.2}\selectfont\renewcommand{\arraystretch}{1.06}
\setlength{\tabcolsep}{10pt}
\begin{tabular}{rcccc}
\toprule
\rowcolor{Pale}State & $\blank$ & $+$ & $0$ & $1$\\
\midrule
0 & $\triple{10}{\blank}{S}$ & $\triple{1}{+}{L}$ & $\triple{0}{0}{R}$ & $\triple{0}{1}{R}$ \\
1 & $\triple{8}{+}{U}$ & $\triple{10}{+}{S}$ & $\triple{1}{0}{D}$ & $\triple{1}{1}{D}$ \\
2 & $\triple{3}{\blank}{L}$ & $\triple{2}{+}{R}$ & $\triple{2}{0}{R}$ & $\triple{2}{1}{R}$ \\
3 & $\triple{10}{\blank}{S}$ & $\triple{9}{\blank}{D}$ & $\triple{4}{\blank}{D}$ & $\triple{5}{\blank}{D}$ \\
4 & $\triple{4}{\blank}{L}$ & $\triple{4}{+}{U}$ & $\triple{7}{0}{D}$ & $\triple{7}{1}{D}$ \\
5 & $\triple{5}{\blank}{L}$ & $\triple{5}{+}{U}$ & $\triple{7}{1}{D}$ & $\triple{6}{0}{L}$ \\
6 & $\triple{7}{1}{D}$ & $\triple{10}{+}{S}$ & $\triple{7}{1}{D}$ & $\triple{6}{0}{L}$ \\
7 & $\triple{7}{\blank}{R}$ & $\triple{1}{\blank}{L}$ & $\triple{10}{0}{S}$ & $\triple{10}{1}{S}$ \\
8 & $\triple{2}{0}{R}$ & $\triple{10}{+}{S}$ & $\triple{2}{0}{R}$ & $\triple{2}{1}{R}$ \\
9 & $\triple{9}{\blank}{L}$ & $\triple{9}{\blank}{U}$ & $\triple{10}{0}{S}$ & $\triple{10}{1}{S}$ \\
10 & $\triple{10}{\blank}{S}$ & $\triple{10}{+}{S}$ & $\triple{10}{0}{S}$ & $\triple{10}{1}{S}$ \\
\bottomrule
\end{tabular}
\end{center}

\vspace{5pt}
\textbf{Metric-Extracted Controller.}\quad Initial state $0$; halt $21$.
\begin{center}
\fontsize{9.4}{11.2}\selectfont\renewcommand{\arraystretch}{1.06}
\setlength{\tabcolsep}{10pt}
\begin{tabular}{rcccc}
\toprule
\rowcolor{Pale}State & $\blank$ & $+$ & $0$ & $1$\\
\midrule
0 & \textcolor{Muted}{--} & \textcolor{Muted}{--} & $\triple{1}{0}{R}$ & $\triple{1}{1}{R}$ \\
1 & \textcolor{Muted}{--} & $\triple{2}{+}{L}$ & $\triple{1}{0}{R}$ & $\triple{1}{1}{R}$ \\
2 & \textcolor{Muted}{--} & \textcolor{Muted}{--} & $\triple{3}{0}{D}$ & $\triple{3}{1}{D}$ \\
3 & $\triple{4}{+}{U}$ & \textcolor{Muted}{--} & \textcolor{Muted}{--} & \textcolor{Muted}{--} \\
4 & $\triple{15}{0}{R}$ & \textcolor{Muted}{--} & $\triple{5}{0}{R}$ & $\triple{5}{1}{R}$ \\
5 & \textcolor{Muted}{--} & $\triple{6}{+}{R}$ & $\triple{5}{0}{R}$ & $\triple{5}{1}{R}$ \\
6 & $\triple{7}{\blank}{L}$ & \textcolor{Muted}{--} & $\triple{6}{0}{R}$ & $\triple{6}{1}{R}$ \\
7 & \textcolor{Muted}{--} & $\triple{16}{\blank}{D}$ & $\triple{12}{\blank}{D}$ & $\triple{8}{\blank}{D}$ \\
8 & $\triple{9}{\blank}{L}$ & \textcolor{Muted}{--} & \textcolor{Muted}{--} & \textcolor{Muted}{--} \\
9 & $\triple{9}{\blank}{L}$ & $\triple{10}{+}{U}$ & \textcolor{Muted}{--} & \textcolor{Muted}{--} \\
10 & \textcolor{Muted}{--} & \textcolor{Muted}{--} & $\triple{11}{1}{D}$ & $\triple{18}{0}{L}$ \\
11 & \textcolor{Muted}{--} & $\triple{3}{\blank}{L}$ & \textcolor{Muted}{--} & \textcolor{Muted}{--} \\
12 & $\triple{13}{\blank}{L}$ & \textcolor{Muted}{--} & \textcolor{Muted}{--} & \textcolor{Muted}{--} \\
13 & $\triple{13}{\blank}{L}$ & $\triple{14}{+}{U}$ & \textcolor{Muted}{--} & \textcolor{Muted}{--} \\
14 & \textcolor{Muted}{--} & \textcolor{Muted}{--} & $\triple{11}{0}{D}$ & $\triple{11}{1}{D}$ \\
15 & \textcolor{Muted}{--} & \textcolor{Muted}{--} & $\triple{5}{0}{R}$ & $\triple{5}{1}{R}$ \\
16 & $\triple{16}{\blank}{L}$ & $\triple{17}{\blank}{U}$ & \textcolor{Muted}{--} & \textcolor{Muted}{--} \\
17 & \textcolor{Muted}{--} & \textcolor{Muted}{--} & $\triple{21}{0}{S}$ & $\triple{21}{1}{S}$ \\
18 & $\triple{20}{1}{D}$ & \textcolor{Muted}{--} & $\triple{19}{1}{D}$ & $\triple{18}{0}{L}$ \\
19 & $\triple{20}{\blank}{R}$ & \textcolor{Muted}{--} & \textcolor{Muted}{--} & \textcolor{Muted}{--} \\
20 & $\triple{20}{\blank}{R}$ & $\triple{3}{\blank}{L}$ & \textcolor{Muted}{--} & \textcolor{Muted}{--} \\
21 & $\triple{21}{\blank}{S}$ & $\triple{21}{+}{S}$ & $\triple{21}{0}{S}$ & $\triple{21}{1}{S}$ \\
\bottomrule
\end{tabular}
\end{center}

\clearpage
\sectiontitle{5 \enspace EQUIVALENCE PROOF}{Equivalence Proof}\label{sec:common5}
\label{sec:part4}\label{sec:equivalence}
\begin{proof}[Proof of Theorem~\ref{thm:main}]
Initially both configurations have the same grid and head,
with control states $0$ related by phase $0$. Suppose the relation holds
before a controller step. Lemma~\ref{lem:protocol} supplies an allowed read symbol
$x$. Equal grids and heads give the same read symbol to both controllers.
Equation~(2) gives equal writes and motions and preserves the relation
between successor states. The next grids and heads therefore agree.

Induction establishes equality of every finite execution prefix. The terminal
condition permits a halt exactly on an edge to the terminal phase.
Lemma~\ref{lem:termination} establishes termination.
Both controllers therefore halt together with equal complete PTOs and final
configurations.
\end{proof}

\textbf{Defined Transitions.}
The metric controller splits expert states according to position in the
algorithm. For example, expert state $4$ (transport and add zero) corresponds
to controller states $12,13,14$, while expert state $9$ (cleanup) corresponds
to $16,17$. The phase invariant excludes every missing entry on valid inputs.
The certificate checks all $43$ defined nonterminal entries. The remaining
four entries contain halt loops. The theorem requires no completion of
undefined transitions.

\bigskip
\sectiontitle{6 \enspace VERIFICATION AND ARTIFACTS}{Verification and Artifacts}\label{sec:common6}
\textbf{Verification.}
Run \code{python3 verify\_equivalence.py {-}{-}quick} from this proof directory to verify
the source fingerprints, equality of the CSV and JSON controller exports,
equality of the exported expert with \code{add\_grid(2)}, and every local
certificate obligation. Omitting \code{{-}{-}quick} also compares two independent
grid executions on 15,876 exhaustive operand pairs (all lengths
$1$--$6$), 44 directed cases up to $1024$ digits, and 128
fixed-seed random cases. All 4,766,713 compared transitions, final grids and
head positions agree. These finite tests supplement the universal proof;
Lemmas~\ref{lem:invariant}--\ref{lem:protocol} and equation~(2)
establish the length-independent claim.

\textbf{Artifact Identity.}
The title page gives the repository-relative case directory. This proof directory
contains the controller and expert exports, reference code, protocol, checker,
and saved verification report. The verification script stores or reads SHA-256 checksums for these artifacts. A checksum identifies the exact bytes
of a file; the checker compares the stored and computed values to detect changes.
These checks bind the certificate to the supplied exports. The invariants,
protocol coverage, and local transition identity establish equivalence.
\end{document}
